\documentclass[12pt,a4paper]{article}
\usepackage{graphicx}
\usepackage{amsmath}
\usepackage{amssymb}
\usepackage{amsthm}
\usepackage{bm}
\usepackage{hyperref}
\usepackage{geometry}
\geometry{margin=2.5cm}

\newtheorem{theorem}{Theorem}

\title{\textbf{CDFD Runtime Paper 03: Dynamics, Regimes, and Stability Tests}}
\author{Steve Bico Mujjabi, MD\\
\small Independent Researcher\\
\small ORCID: \href{https://orcid.org/0009-0001-0556-5516}{0009-0001-0556-5516}}
\date{May 2026}

\begin{document}
\maketitle

\begin{abstract}
\noindent \textbf{Executive synthesis.} This paper states the dynamical core of
CDFD as an executable model. Flow $\Phi$ diffuses through adaptive constraint
$C$; the medium changes through stress response, relaxation, and diffusion; and
the surface-memory product $S M_s$ can stabilize or destabilize the operating
ratio $\Psi_s=(\Phi/C)S M_s$. The local experiments show why bounded numerical
diagnostics matter: the same equations that produce recovery under moderate
stress can produce runaway hyper-adaptation when recovery is too weak.
\end{abstract}

\begin{figure}[ht]
\centering
\fbox{\begin{minipage}{0.92\textwidth}
\centering
\textbf{Runtime evidence path}\\[0.4em]
Prior CDFD mathematics $\longrightarrow$ CDFL/runtime state $\longrightarrow$ bounded output $\longrightarrow$ falsification target.
\end{minipage}}
\caption{Conceptual schematic for the runtime papers. The engine translates the mathematical frame from the prior CDFD papers into executable CDFL/runtime diagnostics; candidate outputs remain tied to explicit tests before any stronger claim is made.}
\end{figure}

\section{Prior CDFD Basis}
This manuscript does not rederive the mathematical core of CDFD. It uses the
Mujjabi Adaptive Operating Ratio and the constrained-vacuum notation introduced
in Part I \cite{MujjabiI2026}, the torus-knot hierarchy and boundary-mode work
from the Part I sequence \cite{MujjabiVI2026,MujjabiIX2026}, the public balance
law developed in the Part I capstone \cite{MujjabiX2026}, the Part I transport
tests \cite{MujjabiXII2026}, and the Life Number and tri-regime biology bridge
from Part II \cite{MujjabiPartII2026}. The runtime papers therefore act as an
execution layer: they keep the mathematics connected to commands, outputs,
falsification tests, and domain-specific limits.


\section{Master Flow Equation}
CDFD begins with the assumption that effective conductivity is inversely related
to constraint. In a coarse spatial model,
\begin{equation}
    \frac{\partial \Phi}{\partial t}
    =
    J_{\mathrm{in}}
    + \nabla \cdot \left(\frac{1}{C}\nabla \Phi\right)
    + S_{\mathrm{src}}
    - D_{\mathrm{sink}},
\end{equation}
where $J_{\mathrm{in}}$ is an optional external input, $S_{\mathrm{src}}$ is a
source term, and $D_{\mathrm{sink}}$ is a dissipation or removal term.

\section{Adaptive Constraint Equation}
The constraint field evolves according to the stress induced by changing flow:
\begin{equation}
    \frac{\partial C}{\partial t}
    =
    \alpha \left|\frac{\partial \Phi}{\partial t}\right|
    - \beta C
    + \gamma \nabla^2 C.
\end{equation}
Here $\alpha$ is stress-driven constraint accumulation, $\beta$ is recovery or
forgetting, and $\gamma$ is spatial spreading of constraint. These parameters
are model parameters, not universal constants.

\section{Surface and Memory Dynamics}
The current engine includes an adaptive surface channel:
\begin{equation}
    \frac{\partial S}{\partial t}
    =
    a\Phi - bC + cM_s - dS + \xi\nabla^2 S.
\end{equation}
The term $\xi\nabla^2 S$ matters because it lets neighboring cells buffer local
surface amplification. Without sufficient recovery or diffusion, $S$ can become
an amplifier rather than a stabilizer.

The structural memory channel is modeled as a hysteresis process:
\begin{equation}
    M_s(t+\Delta t)
    =
    1 + (M_s(t)-1)e^{-\Delta t/\tau_M}
    + \mu H_{\mathrm{overload}}
    + \lambda_M \nabla^2 M_s,
\end{equation}
where $H_{\mathrm{overload}}$ is a bounded overload contribution. This is the
mechanism tested by the vacuum-hysteresis experiment.

\section{Regime Classification}
The runtime uses the operating ratio
\begin{equation}
    \Psi_s = \frac{\Phi}{C}S M_s
\end{equation}
as a local regime variable. Typical labels are:
\begin{itemize}
    \item $\Psi_s < 0.8$: constrained or under-driven;
    \item $0.8 \le \Psi_s \le 1.2$: balanced operating band;
    \item $\Psi_s > 1.2$: overload or stress regime;
    \item sustained high $\Psi_s$: rupture, collapse, escape, or rerouting,
    depending on domain mapping.
\end{itemize}

\section{Experiment Links}
\subsection{Hyper-adaptation boundary}
\texttt{experiments/notebooks/discovery\_adaptive\_surface.py} and the release
selection report identify an overloaded run in which $\Psi_s$ moves from
approximately 5.15 to approximately 979. This output is not a successful engineering result. It is an instability boundary: a test case
showing that high $\Phi$, low $C$, high $S$, high $M_s$, and weak recovery can
produce runaway amplification.

\subsection{Vacuum-memory residual}
\texttt{experiments/notebooks/discovery\_vacuum\_hysteresis.py} tests whether a
localized high-flux pulse leaves a spatially localized $M_s$ residual. The
selected report records peak center $M_s=15.5843$ in
\texttt{experiments/outputs/discovery\_vacuum\_hysteresis.h5}. This is a
candidate simulation result and a future measurement protocol, not empirical
evidence that physical vacuum memory has already been observed.

\section{Numerical Safeguards}
The main runtime now exposes bounded helpers in \texttt{runtime/diagnostics.py}.
These guard toy diagnostics against non-finite values and keep failure modes
visible. A non-finite result is not a discovery. It is a numerical or modeling
failure unless the boundary condition is explicitly being studied as a
breakdown surface.

\section{Conclusion}
The CDFD dynamics are useful only when their regime changes can be reproduced
and falsified. The same mathematical language can describe recovery, overload,
memory, and rupture, but the interpretation must always be tied to the specific
domain mapping and to bounded numerical evidence.
\begin{thebibliography}{9}
\bibitem{MujjabiI2026}
Steve Bico Mujjabi, MD. \emph{Vortex Stability in a Constrained Transport Vacuum and the Origin of the Fine-Structure Constant}. CDFD Part I Paper I, preprint, 2026.

\bibitem{MujjabiVI2026}
Steve Bico Mujjabi, MD. \emph{The Universal Torus Knot Hierarchy: $Z_n$ Symmetry, the Cinquefoil Family, and the General Structure of CDFT Particle Families}. CDFD Part I Paper VI, preprint, 2026.

\bibitem{MujjabiIX2026}
Steve Bico Mujjabi, MD. \emph{Even-$n$ Torus Knots in CDFT: The Exclusion of $T(2,2)$, the Extension to $T(2,2m)$, and the Anti-Phase Mass Scaling Law}. CDFD Part I Paper IX, preprint, 2026.

\bibitem{MujjabiX2026}
Steve Bico Mujjabi, MD. \emph{The Vacuum Equation of State: Deriving Torus Knot Self-Consistency from the Public CDFD Balance Law $\Psi = \Phi/C$}. CDFD Part I Paper X, preprint, 2026.

\bibitem{MujjabiXII2026}
Steve Bico Mujjabi, MD. \emph{Friction, Superconductivity, Plasma States, and Transport Mysteries: Observable Tests of Adaptive Constraint}. CDFD Part I Paper XII, preprint, 2026.

\bibitem{MujjabiPartII2026}
Steve Bico Mujjabi, MD. \emph{CDFD Part II: Origins of Life and Tri-Regime Bioenergetics}. Public release archive, 2026, \url{https://doi.org/10.5281/zenodo.20250821}.
\end{thebibliography}
\end{document}
